% Section 2 -- Related work.
% CITATION POLICY for this section: every \cite key below resolves to an entry in refs.bib. Entries carried over from
% the author's existing bibliography are treated as verified; entries added for this paper are marked in refs.bib as
% needing verification of venue/year/authors before submission. Anything the author knows to exist but that could not
% be cited from a verified record is marked with \open{} rather than invented. See the list at the end of this file.
\section{Related Work}
\label{sec:related}

\subsection{Probing internal representations}

The premise that high-level properties of a model's input are linearly recoverable from its activations predates its
safety application. Structural probes recover syntax from contextual embeddings \cite{hewitt2019structural}, and the
methodological literature on probing classifiers set out early what a probe result does and does not license
\cite{belinkov2022probing}---in particular that a probe's success measures decodability, not use. That distinction is
the ancestor of this paper's concerns, and the reason we insist on baselines and controls rather than on AUC alone.

For properties closer to safety, models have been shown to encode calibration-like signals about their own outputs
\cite{kadavath2022language} and to represent the truth or falsity of statements linearly enough to read off
\cite{azaria2023internal, marks2023geometry}. Representation engineering generalised this into a programme for reading
and steering high-level concepts \cite{zou2023representation}, and activation addition demonstrated behavioural
control from directions obtained without optimisation \cite{turner2023activation}. Related interpretability work
located and edited factual associations directly in weights \cite{meng2022locating}.

\subsection{Activation signals used for safety}

Closest to the present work, SafeSwitch monitors internal activations to detect unsafe generations and steer away from
them \cite{han2025safeswitch}, and ConceptGuard builds guardrails from interpretable jailbreak concepts
\cite{conceptguard2025}. Both share this paper's premise---that a signal already present in the forward pass can be
used for safety at negligible cost---and both raise the evaluation question this paper is about: what has to be
controlled before such a number can be believed.

Three further lines bound what a probe result can be taken to mean. Probes have been shown to be sensitive to
evaluation awareness and to input format, and at least one study length-matches its conditions for exactly the reason
we do \cite{evalawareness2026}. Monitors degrade as the distribution they were fitted on ages
\cite{staleness2026}, which is the temporal counterpart of the extraction-convention effects reported here. And an
attack literature constructs inputs whose activations evade probe detection while the behaviour persists
\cite{bailey2025obfuscated}, which is the
adversarial counterpart of our distributional negative result in Section~\ref{sec:obfuscation}: we show a probe fails
on encodings chosen without reference to it, and that line shows it fails on encodings chosen against it. Work on detecting
high-stakes interactions with activation probes \cite{mckenzie2025probes} sits alongside both, is the closest
published analogue to the minimal-data question we ask, and already compares token-aggregation schemes for
probes---last-token, mean, max and attention-weighted pooling---including in small-data settings. Aggregation as such
is therefore not what Section~\ref{sec:read} contributes. What is new there is localizing the read to the
threat-bearing span of an agent scaffold rather than pooling over the input as a whole---measured against that
whole-input baseline, localisation is worth $+0.02$ to $+0.09$ AUC on the models where the probe works
(Table~\ref{tab:aag_designs})---and bounding what the deployable reads achieve against a label-informed oracle placed
at the injection itself.

\subsection{Probes in production}
\label{sec:production}

The minimal-data question is worth asking because the maximal-data version has already been answered affirmatively at
frontier scale. Two January 2026 reports describe probe-based monitoring in production: Constitutional
Classifiers++ at Anthropic \cite{cunningham2026classifiers} and a production probe stack for Gemini at Google
DeepMind \cite{gdm2026probes}. Read together they converge on an architecture---probe-first cascades, in which a cheap
activation-space detector triages and a more expensive classifier adjudicates what it flags---while differing in
nearly every implementation choice: all-layers versus middle-layer reads, exchange-level versus input-level scoring,
and CBRN versus cyber threat models. That two independent teams converged on the shape while diverging on the details
is the practical case for treating those details as first-class experimental variables, which is what this paper does.

One result in the DeepMind report bears directly on our thesis. Their probes degrade under a short-to-long context
distribution shift: a detector fitted on short inputs does not hold up when the input it must score is long
\cite{gdm2026probes}. That is the frontier-scale form of the observation in Section~\ref{sec:read}---that the text a
probe is read from, and where in it the read happens, can matter more than how the probe was trained---and it is
evidence that extraction convention is not an artifact of the minimal-data regime we study.

\subsection{Guard models}

The deployed alternative to a probe is a separate classifier. Llama Guard established the pattern of an
instruction-tuned model applying an explicit hazard taxonomy to inputs and outputs \cite{inan2023llama}; ShieldGemma
\cite{zeng2024shieldgemma} and WildGuard \cite{han2024wildguard} extend it with different taxonomies and training
mixtures. Section~\ref{sec:guard} compares against Llama Guard~3 directly and finds it more accurate than the probe on
every metric, which is why this paper's positive claims are about cost structure and about evaluation methodology
rather than about beating a guard model.

\subsection{Prompt injection}

Indirect prompt injection---instructions smuggled into content the model retrieves rather than into the user's
turn---was characterised as a systemic risk for tool-integrated applications by Greshake et al.
\cite{greshake2023youve}, formalised and benchmarked by Liu et al. \cite{liu2024formalizing}, and given standard
evaluation harnesses by InjecAgent \cite{zhan2024injecagent} and AgentDojo \cite{debenedetti2024agentdojo}. Defences
divide broadly into training-time approaches that teach the model to privilege the user's instruction
\cite{wallace2024instruction}, architectural approaches that constrain what injected text can reach
\cite{debenedetti2025camel}, and detection approaches, of which the probe studied in Section~\ref{sec:read} is one.
Our contribution to this line is not a new defence but an observation about how the detection variety is measured: the
signal's position inside the scaffold matters more than the layer, and the standard benchmark's small pool of
injection strings makes a case-level split measure recognition rather than generalisation.

\subsection{Benchmarks and their validity}

We use HarmBench \cite{mazeika2024harmbench} for harmful behaviours, XSTest \cite{rottger2024xstest} for
exaggerated-safety false positives, JailbreakBench \cite{chao2024jailbreakbench} for on-topic benign twins,
InjecAgent \cite{zhan2024injecagent} and AgentDojo \cite{debenedetti2024agentdojo} for agent injection, and AILuminate
\cite{mlcommons2025ailuminate} for the obfuscation study. XSTest and JailbreakBench are used together deliberately:
Section~\ref{sec:data} explains why a single pooled benign set conceals the failure mode this paper is most concerned
with.

The methodological findings in this paper belong to an older tradition than the probing literature. Construct
validity---whether a measurement instrument measures the construct it names, rather than something correlated with
it---is the standard frame for our Section~\ref{sec:learn} result, and the classical statement of it is Messick's
\cite{messick1995validity}. The critique that benchmarks are routinely treated as valid measures of
capacities they were never shown to measure has been made directly of machine-learning benchmarks
\cite{raji2021benchmark}, and a recent audit applies it to agent-safety benchmarks specifically
\cite{agentsafetyaudit2026}. Our contribution to this line is narrow and empirical: we exhibit an evaluation set that
cannot certify the semantic claim attached to it, and we show what it costs to find that out.

\subsection{Positioning}

The three probe families studied here---harmful-content filtering, agent-injection detection and policy compliance,
named \af{}, \aag{} and \apc{} in Section~\ref{sec:tasks}---and the paired contrastive data construction examined in
Section~\ref{sec:learn}, originate in a publicly released minimal-data recipe \cite{messenger2026aase}; this paper's
contribution relative to it is diagnosis and controls---extraction conventions and evaluation validity---rather than a
new probe family.

The papers above establish that safety-relevant properties are decodable from activations and that probes built on
that fact can be made to work. This paper takes the next question seriously: under what evaluation conditions is a
reported probe number evidence of the claim attached to it. Its two positive results are repairs to extraction
conventions rather than new probe architectures, and its negative results concern what minimal-data probes turn out to
have learned when the obvious confounds are controlled.

% =====================================================================================================
% LITERATURE SEARCH -- CLOSED 2026-09-05. The author's search resolved all six open items; citations are in
% refs.bib. Two classes of residual work remain, both proof-stage and neither affecting this text:
%   (a) entries marked NEEDS-TITLE / NEEDS-AUTHORS / NEEDS-URL in refs.bib, where the search supplied an
%       identifier but not a complete record. The \cite keys are stable, so completing them will not touch
%       any section file.
%   (b) RESOLVED 2026-09-05: the three [verify] citations were confirmed by the author (Bailey et al. ICML 2025,
%       Messick 1995 Am. Psychol. 50(9) 741-749, Raji et al. NeurIPS 2021 D&B). Only Bailey's final venue string
%       needs a look at the arXiv page at proof stage.
% =====================================================================================================
